%------------------------------------------------------------------------------%
\begin{question}
\begin{columns}
\column{0.5\linewidth}
  Encontre o vetor tangente ao astroide
  \[
    x = \cos^3(t) \qquad
  \]

  \[
    y = \sin^3(t) \qquad
  \]

  \[
    0 \leq t \leq 2\pi
  \]

  no ponto \(t=\frac{\pi}{6}\)
\column{0.5\linewidth}
  \pause
  ~\\[5mm]
  \includegraphics{B_04_exemplo_5_curva}
\end{columns}
\end{question}

%------------------------------------------------------------------------------%
\begin{resolution}{Solução}
  Avaliando das derivadas das componentes

  \pause
  \medskip
	\[
		\frac{dx}{dt}
		\pause
		\;=\; \frac{d}{dt}\cos^3(t)
		\pause
		\;=\; 3\cos^2(t)\frac{d}{dt}\cos(t)
		\pause
		\;=\; -3\cos^2(t)\sin(t)
	\]

	\pause
  \medskip
	\[
		\frac{dy}{dt}
		\pause
		\;=\; \frac{d}{dt}\sin^3(t)
		\pause
		\;=\; 3\sin^2(t)\frac{d}{dt}\sin(t)
		\pause
		\;=\; 3\sin^2(t)\cos(t)
	\]
\end{question}

%------------------------------------------------------------------------------%
\begin{resolution}{Solução}
  Avaliando a derivada de $x$ no ponto \(t=\frac{\pi}{6}\)
  \pause
  \[
    \frac{dx}{dt}\left(\frac{\pi}{6}\right)
    \pause
    \;=\; \left[ -3\cos^2(t)\sin(t)\right] \evalat{\frac{\pi}{6}}{}
    \pause
    \;=\; -3\cos^2\left(\frac{\pi}{6}\right)\sin\left(\frac{\pi}{6}\right)
  \]
  \[
    \pause
    \hphantom{\frac{dx}{dt}\left(\frac{\pi}{6}\right)}
    \;=\; -3\left(\frac{\sqrt{3}}{2}\right)^2 \left(\frac{1}{2}\right)
    \pause
    \;=\; -3\; \frac{3}{4}\; \frac{1}{2}
    \pause
    \;=\; \frac{-9}{8}
  \]
\end{resolution}

%------------------------------------------------------------------------------%
\begin{resolution}{Solução}
  Avaliando a derivada de $y$ no ponto \(t=\frac{\pi}{6}\)
  \pause
  \[
    \frac{dy}{dt}\left(\frac{\pi}{6}\right)
    \pause
    \;=\; \left[ 3\sin^2(t)\cos(t)\right] \evalat{\frac{\pi}{6}}{}
    \pause
    \;=\; 3\sin^2\left(\frac{\pi}{6}\right)\cos\left(\frac{\pi}{6}\right)
  \]
  \[
    \pause
    \hphantom{\frac{dy}{dt}\left(\frac{\pi}{6}\right)}
    \;=\; 3\left(\frac{1}{2}\right)^2 \left(\frac{\sqrt{3}}{2}\right)
    \pause
    \;=\; 3 \; \frac{1}{4} \; \frac{\sqrt{3}}{2}
    \pause
    \;=\; \frac{3\sqrt{3}}{8}
  \]
\end{resolution}

%------------------------------------------------------------------------------%
\begin{resolution}{Solução}
  Função vetor tangente
  \pause
  \[
    v(t)
    \pause
    \;=\; \frac{d}{dt}\vetor{x}{y}
    \pause
    \;=\; \vetor{-3\cos^2(t)\sin(t)}{\hpm 3\sin^2(t)\cos(t)}
  \]

  \pause
  No ponto \(t=\frac{\pi}{6}\)
  \[
    \pause
    v\left(\frac{\pi}{6}\right)
    \pause
    \;=\; \Vetor{\frac{-9}{8}}{\frac{3\sqrt{3}}{8}}
    \pause
    \;=\; \frac{3}{8} \vetor{-3}{\sqrt{3}}
  \]
\end{resolution}

%------------------------------------------------------------------------------%
\begin{resolution}{Solução}
  \centering
  \includegraphics{B_04_exemplo_5_vetor}
\end{resolution}

%------------------------------------------------------------------------------%
